\documentclass{article}
\usepackage{amsmath,amsthm}
\title{Compound Locators in Spectral Perturbation}
\author{Synthetic Case P5-05}
\newtheorem{thm}{Theorem}[section]
\newtheorem{cor}[thm]{Corollary}
\begin{document}
\maketitle

\section{Perturbation Bounds}
\label{sec:perturb}

\begin{thm}[Resolvent continuity]
\label{thm:resolvent}
When self-adjoint operators converge in norm resolvent sense, isolated eigenvalues converge continuously.
\end{thm}

\begin{cor}[Simple eigenvalue rate]
\label{cor:simple-rate}
For a simple eigenvalue the shift is linear in the perturbation amplitude at leading order.
\end{cor}

\begin{cor}[Eigenspace projection bound]
\label{cor:proj-bound}
The spectral projection norm error is bounded by the resolvent separation ratio.
\end{cor}

\end{document}
